\subsection{Experiments Added in v2 After Review (Post Hoc)}\label{subsec:v2}

\noindent\textbf{Status.} This subsection reports experiments added in v2 after review. They answer four review findings: no consistency-weight or dwell-plus-improvement selector among the comparators, a noise level that was neither swept nor reported per scenario, no scripted target for the zero-switch outcomes, and no test of progress hardening near a completed maneuver. They are post hoc, because they were run after the v1 results and the review were known. The plan, with the metrics and the tuning procedure, was committed before the first run (\code{paper/arxiv\_ref/V2\_EXPERIMENT\_PLAN.md}). No v1 number changes. The first six rows of \cref{tab:v2native} are the v1 rows, regenerated by the new driver and checked equal to \code{ablation\_raw.csv} row by row. The harness, the five scenarios, seeds 0--49 and the observer are the v1 ones. All results are libstdc++ results, and no libc++ rerun was made. The raw data are in \code{src/compass\_eval/results/v2\_experiments/}; \code{scripts/test\_v2\_experiments.sh} rebuilds them and compares every file byte for byte.

\noindent\textbf{Scripted targets.} Each scenario gets a target side from its script alone, by one rule: the side that the scripted mean advantage favours at the end of the horizon, with a short spike not counted. The start side is R. The targets are R in \code{transient\_spike}, L in \code{mid\_reversal}, \code{clean\_commit} and \code{intermittent}, and none in \code{near\_tie}. A target is a property of the script, not an independent label, and it says nothing about response time. \code{mid\_reversal} gets a target under the same rule as the other scenarios: its advantage is $-0.25$ at the start and $+0.25$ at the end. The ``correct final side'' column counts the runs whose last decision equals the target.

\noindent\textbf{Two further selectors.} Neither reproduces a published method. CW, in the spirit of \citet{degroot2024topology} and SHINE~\citep{martinezbaselga2026shine}, adds a bonus $b$ for the previous class to the cost comparison. It is an instantaneous margin without memory. The bonus is additive on the previous class, whereas the T-MPC weight of \citet{degroot2024topology} is multiplicative on the previous planner's cost, so CW is a simplified analogue. LI, in the spirit of \citet{li2026bigcbf}, switches only when the time since the last switch is at least $T_{\min}$ and the best class improves on the current one by at least a fraction $r$ of its cost. It has no accumulation. Both act on the same per-cycle cost streams as the other rules. The comparison is therefore between selection rules on scripted label streams, not between the published planners.

\noindent\textbf{One tuning procedure for every rule.} Each rule with a free parameter is tuned by the same procedure, fixed before the runs (\cref{tab:v2tuning}). \sname\ gets the same treatment: its $E_0$ is tuned, and the default $E_0=0.30$ is a grid point. The selected value of CW ($b=0.4$) and of LI lies at the edge of its grid, and the grids were not extended. Tuning and evaluation use the same scenario scripts with different seeds, so every tuned row, \sname\ included, is favoured over a rule fixed in advance. The full grids, with tuning and evaluation losses, are in \code{v2\_experiments/V2\_results.md}.

\begin{table}[t]
\centering
\caption{\textbf{Added in v2 after review (post hoc): tuning procedure and selected values.} One procedure for all rules: grid search on tuning seeds 1000--1049 (disjoint from the evaluation seeds), loss = excess switches $+\,3\times$ [wrong final side] averaged over the five scenarios at native noise; ties go to the earlier grid point. The selected point is applied unchanged to the evaluation seeds. Evaluation loss uses seeds 0--49. Selected values at the edge of a grid are not extended.}
\label{tab:v2tuning}
\small
\resizebox{\linewidth}{!}{%
\begin{tabular}{lllcc}
\toprule
Rule & Grid & Selected & Tuning loss & Evaluation loss \\
\midrule
COMPASS $E_0$ & \{0.10, 0.15, 0.20, 0.30, 0.40\} & 0.15 & 0.024 & 0.048 \\
Simple dwell $T_{\mathrm{dwell}}$ (s) & \{0.2, 0.4, 0.6, 1.0, 1.5, 2.0\} & 0.4 & 0.468 & 0.476 \\
CW bonus $b$ & \{0, 0.05, 0.10, 0.15, 0.20, 0.30, 0.40\} & 0.4 & 0.628 & 0.608 \\
LI $(T_{\min}$ s, $r)$ & $\{0.3, 0.6, 1, 2\}\times\{0.05, 0.10, 0.20, 0.40\}$ & 2, 0.4 & 1.024 & 1.056 \\
\bottomrule
\end{tabular}%
}
\end{table}

\begin{table}[t]
\centering
\caption{\textbf{Added in v2 after review (post hoc): comparison with two further selection rules and tuned variants, native noise.} Mean$\pm$SD over $N=250$ runs per row (5 scenarios $\times$ 50 seeds, seeds 0--49), $\sigma=0.12$ in \texttt{near\_tie} and $0.03$ elsewhere. Censored = number of the 250 runs without a reported $t_{\text{sfx}}$. Correct final side = runs whose last decision equals the scripted target side (\cref{tab:v2native_scn}), pooled over the four scenarios with a target (out of 200). The first six rows repeat the archived rows of \cref{tab:main} unchanged; no column is marked best. Tuned rows use the single procedure of \cref{tab:v2tuning}.}
\label{tab:v2native}
\small
\resizebox{\linewidth}{!}{%
\begin{tabular}{lcccccc}
\toprule
Rule & Switches & Sign changes (/s) & Entropy (bits) & $t_{\text{sfx}}$ (s) & Censored & Correct final side \\
\midrule
Full (proposed) & 0.40$\pm$0.49 & 0.040$\pm$0.049 & 0.018$\pm$0.022 & 1.33$\pm$1.66 & 0 & 150 \\
$-$Hysteresis & 0.40$\pm$0.49 & 0.040$\pm$0.049 & 0.018$\pm$0.022 & 1.00$\pm$1.20 & 0 & 150 \\
$-$Progress hardening & 0.40$\pm$0.49 & 0.040$\pm$0.049 & 0.018$\pm$0.022 & 1.31$\pm$1.65 & 0 & 150 \\
$-$Accumulator (argmin) & 29.74$\pm$36.70 & 2.974$\pm$3.670 & 0.420$\pm$0.374 & 3.60$\pm$4.64 & 72 & 200 \\
Simple dwell (0.6 s) & 0.80$\pm$0.75 & 0.080$\pm$0.075 & 0.034$\pm$0.031 & 2.26$\pm$3.90 & 50 & 150 \\
$-$Class correspondence & 37.20$\pm$36.12 & 3.720$\pm$3.612 & 0.515$\pm$0.401 & 3.60$\pm$4.63 & 71 & 200 \\
COMPASS, $\lambda=1$ & 0.60$\pm$0.49 & 0.060$\pm$0.049 & 0.027$\pm$0.022 & 2.81$\pm$3.71 & 50 & 200 \\
COMPASS, tuned $E_0=0.15$ & 0.65$\pm$0.53 & 0.065$\pm$0.053 & 0.029$\pm$0.023 & 2.49$\pm$3.80 & 50 & 200 \\
Simple dwell, tuned $T_{\mathrm{dwell}}=0.4$ s & 1.08$\pm$0.57 & 0.108$\pm$0.057 & 0.047$\pm$0.023 & 2.82$\pm$4.09 & 58 & 200 \\
CW, tuned $b=0.4$ & 0.11$\pm$0.32 & 0.011$\pm$0.032 & 0.005$\pm$0.014 & 0.14$\pm$0.82 & 1 & 150 \\
LI, tuned $T_{\min}=2$ s, $r=0.4$ & 0.78$\pm$0.82 & 0.078$\pm$0.082 & 0.033$\pm$0.032 & 4.13$\pm$4.40 & 58 & 160 \\
\bottomrule
\end{tabular}%
}
\end{table}

\begin{table}[t]
\centering
\caption{\textbf{Added in v2 after review (post hoc): per scenario, mean switches per encounter; runs with the scripted final side (of 50).} Cell = mean switches; correct final side. Targets (scripted, \cref{subsec:v2}): tie none (marked --), spike R, rev L, clean L, interm.\ L. Switch counts exclude the initial class to first output, so a rule that commits on its first output shows 0 switches and is scored by the final-side count. Scenario names: tie = \texttt{near\_tie}, spike = \texttt{transient\_spike}, rev = \texttt{mid\_reversal}, clean = \texttt{clean\_commit}, interm. = \texttt{intermittent}.}
\label{tab:v2native_scn}
\small
\resizebox{\linewidth}{!}{%
\begin{tabular}{lccccc}
\toprule
Rule & tie & spike & rev & clean & interm. \\
\midrule
Full (proposed) & 0.00; -- & 0.00; 50 & 0.00; 0 & 1.00; 50 & 1.00; 50 \\
$-$Hysteresis & 0.00; -- & 0.00; 50 & 0.00; 0 & 1.00; 50 & 1.00; 50 \\
$-$Progress hardening & 0.00; -- & 0.00; 50 & 0.00; 0 & 1.00; 50 & 1.00; 50 \\
$-$Accumulator (argmin) & 98.16; -- & 2.00; 50 & 13.44; 50 & 0.00; 50 & 35.12; 50 \\
Simple dwell (0.6 s) & 0.00; -- & 2.00; 50 & 1.00; 50 & 1.00; 50 & 0.00; 0 \\
$-$Class correspondence & 98.12; -- & 2.00; 50 & 46.00; 50 & 0.00; 50 & 39.88; 50 \\
COMPASS, $\lambda=1$ & 0.00; -- & 0.00; 50 & 1.00; 50 & 1.00; 50 & 1.00; 50 \\
COMPASS, tuned $E_0=0.15$ & 0.00; -- & 0.24; 50 & 1.00; 50 & 1.00; 50 & 1.00; 50 \\
Simple dwell, tuned $T_{\mathrm{dwell}}=0.4$ s & 0.38; -- & 2.00; 50 & 1.00; 50 & 1.00; 50 & 1.00; 50 \\
CW, tuned $b=0.4$ & 0.04; -- & 0.00; 50 & 0.00; 0 & 0.50; 50 & 0.02; 50 \\
LI, tuned $T_{\min}=2$ s, $r=0.4$ & 1.68; -- & 1.20; 10 & 1.00; 50 & 0.00; 50 & 0.00; 50 \\
\bottomrule
\end{tabular}%
}
\end{table}


\noindent\textbf{Results at native noise.} (i) Default \sname\ reaches the scripted target side in none of the 50 \code{mid\_reversal} runs. So do $-$hysteresis and $-$progress hardening. The zero-switch outcome of \cref{subsec:ablation} is therefore a miss against the script's own final-window advantage. Simple dwell (0.6~s) reaches that target in 50 of 50 runs and misses the \code{intermittent} target in 50 of 50 runs. (ii) Tuned CW ($b=0.4$) has 0.11 switches per encounter against 0.40 for default \sname, a mean $t_{\text{sfx}}$ of 0.14~s against 1.33~s, and 150 correct final sides out of 200 for both. It has one censored run, where default \sname\ has none. It also misses the \code{mid\_reversal} target in all 50 runs. A tuned instantaneous margin thus matches or beats the default accumulator on these offline statistics at native noise. The switch counts exclude the initial output (\cref{sec:exp}), so the final-side column is the safer comparison. (iii) Among the tuned rules, \sname\ with $E_0=0.15$ has the lowest evaluation loss (0.048, against 0.476 for dwell, 0.608 for CW and 1.056 for LI; default \sname\ has 0.600). Tuned \sname\ and tuned dwell both reach the target in all 200 runs. Tuned dwell switches 0.38 times per \code{near\_tie} encounter and twice per \code{transient\_spike} encounter; tuned \sname\ switches 0.00 and 0.24 times. (iv) Tuned LI reaches the \code{transient\_spike} target in only 10 of 50 runs. The ordering of tuned rules holds for these five scripts at one noise level. It is not evidence about other scenarios.

\begin{table}[t]
\centering
\caption{\textbf{Added in v2 after review (post hoc): noise sweep, mean switches per encounter.} Cell = mean over the five scenarios (mean in \texttt{near\_tie}), 250 runs per cell (50 in \texttt{near\_tie}). The same $\sigma$ is used in all five scenarios. Tuned rules keep the value tuned at native noise. Reference: the stationary SD of the unclipped leaky accumulation, $\sigma\,dt/\sqrt{1-\lambda^2}$, is $0.006$, $0.012$, $0.025$, $0.041$, $0.062$ for the five $\sigma$ against $E_0=0.30$.}
\label{tab:v2sigma_sw}
\small
\resizebox{\linewidth}{!}{%
\begin{tabular}{lccccc}
\toprule
Rule & $\sigma=0.03$ & 0.06 & 0.12 & 0.20 & 0.30 \\
\midrule
Full (proposed) & 0.40 (0.00) & 0.40 (0.00) & 0.40 (0.00) & 0.40 (0.00) & 0.40 (0.00) \\
$-$Hysteresis & 0.40 (0.00) & 0.40 (0.00) & 0.40 (0.00) & 0.40 (0.00) & 0.47 (0.00) \\
$-$Progress hardening & 0.40 (0.00) & 0.40 (0.00) & 0.40 (0.00) & 0.40 (0.00) & 0.45 (0.00) \\
COMPASS, $\lambda=1$ & 0.60 (0.00) & 0.58 (0.00) & 0.57 (0.00) & 0.57 (0.00) & 0.59 (0.00) \\
COMPASS, tuned $E_0=0.15$ & 0.65 (0.00) & 0.71 (0.00) & 0.74 (0.00) & 0.82 (0.00) & 0.85 (0.00) \\
$-$Accumulator (argmin) & 29.74 (98.16) & 33.73 (98.16) & 40.40 (98.16) & 53.87 (98.16) & 68.11 (98.16) \\
Simple dwell (0.6 s) & 0.80 (0.00) & 0.81 (0.00) & 0.88 (0.00) & 0.68 (0.00) & 0.46 (0.00) \\
Simple dwell, tuned $T_{\mathrm{dwell}}=0.4$ s & 1.08 (0.38) & 1.08 (0.38) & 1.08 (0.38) & 1.00 (0.38) & 0.92 (0.38) \\
CW, tuned $b=0.4$ & 0.10 (0.00) & 0.16 (0.00) & 0.77 (0.04) & 2.99 (5.12) & 10.89 (18.46) \\
LI, tuned $T_{\min}=2$ s, $r=0.4$ & 0.44 (0.00) & 0.61 (0.04) & 1.09 (1.68) & 2.57 (3.98) & 3.92 (4.72) \\
\bottomrule
\end{tabular}%
}
\end{table}

\begin{table}[t]
\centering
\caption{\textbf{Added in v2 after review (post hoc): noise sweep, share of runs with the scripted final side.} Share of the 200 runs (four scenarios with a target $\times$ 50 seeds) whose last decision equals the target side of \cref{tab:v2native_scn}; same runs as \cref{tab:v2sigma_sw}.}
\label{tab:v2sigma_hit}
\small
\resizebox{\linewidth}{!}{%
\begin{tabular}{lccccc}
\toprule
Rule & $\sigma=0.03$ & 0.06 & 0.12 & 0.20 & 0.30 \\
\midrule
Full (proposed) & 0.75 & 0.75 & 0.75 & 0.75 & 0.75 \\
$-$Hysteresis & 0.75 & 0.75 & 0.75 & 0.75 & 0.82 \\
$-$Progress hardening & 0.75 & 0.75 & 0.75 & 0.75 & 0.81 \\
COMPASS, $\lambda=1$ & 1.00 & 0.97 & 0.96 & 0.96 & 0.96 \\
COMPASS, tuned $E_0=0.15$ & 1.00 & 0.99 & 0.96 & 0.96 & 0.97 \\
$-$Accumulator (argmin) & 1.00 & 1.00 & 0.98 & 0.93 & 0.89 \\
Simple dwell (0.6 s) & 0.75 & 0.76 & 0.89 & 0.84 & 0.72 \\
Simple dwell, tuned $T_{\mathrm{dwell}}=0.4$ s & 1.00 & 1.00 & 1.00 & 0.98 & 0.95 \\
CW, tuned $b=0.4$ & 0.75 & 0.74 & 0.98 & 1.00 & 0.97 \\
LI, tuned $T_{\min}=2$ s, $r=0.4$ & 0.80 & 1.00 & 0.99 & 0.83 & 0.65 \\
\bottomrule
\end{tabular}%
}
\end{table}


\noindent\textbf{Noise sweep and the leak.} The sweep uses the same $\sigma$ in all five scenarios (0.03 to 0.30), and the tuned rules keep their native-noise values. (i) Default \sname\ switches 0.00 times per \code{near\_tie} encounter at every $\sigma$. Its pooled mean stays at 0.40 up to $\sigma=0.20$, and so does that of $-$hysteresis and $-$progress hardening. At $\sigma=0.30$ the two ablations change in mixed directions: 0.47 and 0.45 pooled switches against 0.40, and 0.82 and 0.81 of the targets reached against 0.75. (ii) The $\lambda=1$ variant also switches 0.00 times in \code{near\_tie} at every $\sigma$, and reaches 0.96 to 1.00 of the targets. The $\lambda=1$ variant keeps $\delta_{\text{floor}}=0.05$ and $-$hysteresis keeps the leak, so no cell removes both. Removing the leak alone ($\lambda=1$, margin kept) therefore does not increase \code{near\_tie} switching in the swept range. This does not show that threshold accumulation alone suffices. In these scripts the leak shows up as slower response: with $\lambda=1$ the \code{mid\_reversal} target is reached in 50 of 50 runs, and with $\lambda=0.97$ in none. Leak and threshold are thus partly separated, and the v1 wording ``leaky accumulator'' as the suppressor is not supported as stated; the sweeps do not isolate one mechanism. (iii) Tuned CW has 0.04 switches per \code{near\_tie} encounter at $\sigma=0.12$, 5.12 at $\sigma=0.20$ and 18.46 at $\sigma=0.30$. Tuned LI has 1.68 at $\sigma=0.12$ and 4.72 at $\sigma=0.30$. The selected CW bonus $b=0.4$ sits at the edge of the pre-declared grid, which was not extended, so the high-noise CW numbers may overstate CW's degradation. The accumulator's advantage over tuned CW appears from $\sigma=0.20$ on, above the native noise of the v1 scenarios. Over tuned LI it already appears at $\sigma=0.12$. (iv) The \code{near\_tie} counts of immediate argmin (98.16) and of tuned dwell (0.38) do not depend on $\sigma$, because the scripted tie advantage is zero and only the sign of the noise matters. (v) The sweep does not reach the regime where the accumulation noise nears $E_0$. The stationary SD of the unclipped accumulation, $\sigma\,dt/\sqrt{1-\lambda^2}$, is 0.062 at $\sigma=0.30$ against $E_0=0.30$. It would reach $E_0$ near $\sigma=1.46$. Correlated noise was not swept.

\begin{table}[t]
\centering
\caption{\textbf{Added in v2 after review (post hoc): late spurious reversal.} Scripted burst of advantage toward R for $N_b$ cycles from cycle 150, after the maneuver toward L has progressed ($v_{\text{lat}}=0.15$~m/s, $\sigma=0.03$, 50 seeds per cell). Cell = mean switches; number of the 50 runs that left L during or after the burst. $\rho$ = mean progress variable at the start of the burst (core rows only). Target: stay on L (the burst is spurious by construction).}
\label{tab:v2late}
\small
\resizebox{\linewidth}{!}{%
\begin{tabular}{lccccc}
\toprule
Rule & $\rho$ at burst start & $N_b=20$ & 30 & 40 & 60 \\
\midrule
Full (proposed) & 0.94 & 1.00; 0 & 1.00; 0 & 1.00; 0 & 1.00; 0 \\
$-$Progress hardening & 0.95 & 1.00; 0 & 2.00; 50 & 2.00; 50 & 2.00; 50 \\
$-$Hysteresis & 0.98 & 1.00; 0 & 1.00; 0 & 1.00; 0 & 1.00; 0 \\
COMPASS, $\lambda=1$ & 1.00 & 1.00; 0 & 1.00; 0 & 1.00; 0 & 1.00; 0 \\
$-$Accumulator (argmin) & -- & 2.00; 50 & 2.00; 50 & 2.00; 50 & 1.00; 50 \\
Simple dwell (0.6 s) & -- & 3.00; 50 & 3.00; 50 & 2.00; 50 & 2.00; 50 \\
Simple dwell, tuned $T_{\mathrm{dwell}}=0.4$ s & -- & 3.00; 50 & 3.00; 50 & 3.00; 50 & 2.00; 50 \\
CW, tuned $b=0.4$ & -- & 2.50; 50 & 2.50; 50 & 2.50; 50 & 1.50; 50 \\
LI, tuned $T_{\min}=2$ s, $r=0.4$ & -- & 2.00; 50 & 2.00; 50 & 2.00; 50 & 1.00; 50 \\
\bottomrule
\end{tabular}%
}
\end{table}


\noindent\textbf{Late spurious reversal.} The scenario was built to probe progress hardening: advantage $+0.40$ for 150 cycles, then a burst of $-0.40$ for $N_b$ cycles, then $+0.40$ again, with a progress input of $0.15$~m/s. The plan estimated $\rho\approx0.77$ at the start of the burst. The observed mean for default \sname\ is 0.94, because the commitment came earlier than the estimate assumed. (i) Default \sname\ stays on L in all runs for every $N_b$ (0 of 50 leave). So do $-$hysteresis and the $\lambda=1$ variant. This hold follows by construction from the P4 clip lockout once $\rho$ exceeds $\bar\rho$ ($\Eth(\rho)>\Eremax$), so it confirms the implementation rather than measuring robustness. (ii) $-$Progress hardening ($k_\rho=0$) does not leave L for $N_b=20$ (0 of 50) and leaves it in 50 of 50 runs for $N_b=30$, 40 and 60. Progress hardening is therefore the only tested mechanism that holds the maneuver against a burst of 30 cycles or longer. (iii) Immediate argmin, simple dwell, tuned dwell, tuned CW and tuned LI leave L in 50 of 50 runs for every $N_b$. None of them uses progress. For $N_b=60$ the burst lasts until the end of the log. This scenario was designed so that the mechanism could matter. The result shows that it can matter in this script. It does not show that such bursts occur in practice, and it does not show that holding the maneuver is the better decision.

\noindent\textbf{What this changes in the claims.} (i) The comparison with a consistency weight is no longer absent, but it is limited to a self-tuned selector on scripted label streams. It is not a reproduction of the published planners, and opinion-dynamics passing and closed-loop evaluation remain unrun. (ii) At native noise a tuned memoryless margin is as good as the default accumulator on switch count and $t_{\text{sfx}}$; the accumulator's measured advantage is in tie suppression against CW at $\sigma\ge0.20$ and against LI at $\sigma=0.12$, and in fewer unnecessary switches than a tuned dwell timer. (iii) The default \sname\ is slow on \code{mid\_reversal}, and a lower $E_0$ or $\lambda=1$ removes the miss in these scripts. (iv) Hysteresis and progress hardening showed no regression through $\sigma=0.20$ in the sweep; progress hardening matters in the late-burst script. Not shown: closed-loop behaviour, comparison with the published planners, noise comparable to $E_0$, correlated noise, and any human or physical result.

